Our main target of ESO-GAD is to include the relational semantics of an edge into the spectral filtering in the frequency domain.
Recent evidence on edge heterophily discrimination shows
that homophilic and heterophilic edges can be estimated from local features~\citep{GREET, H2-FDetector}.
We therefore cast edge-semantics estimation as a supervised task, which classifies edges as heterophilic or homophilic given the input training data.
For each edge $(u,v)\in\cE$ we learn a heterophily probability
\begin{equation}
  q_{uv} \;=\; \Pr\!\left[\, y_u \neq y_v \,\middle|\, \vz_{uv} \,\right]
  \;\in\; [0,1],
  \label{eq:quv}
\end{equation}
where $\vz_{uv}$ is a symmetric edge descriptor and the probability is parameterized by a small multilayer perceptron (MLP).
The probabilistic design is necessary in ESO-GAD because the labeled set $\cV_L$ is tiny and severely imbalanced.
The output probability $q_{uv}$ can be considered as a smooth transition of an edge semantic between the homophily label and heterophily label.
Unlike hard pruning, where each false positive deletes the local structure, a calibrated soft probability is expected to degrade gracefully when there is uncertainty in classifying the edge.

The edge descriptor $\vz_{uv}$ comprises a feature component $\vz^{x}_{uv}$ and a local structural component $\vz^{s}_{uv}$.
Comparing endpoint features is a standard basis for edge representations in various graph learning tasks ~\citep{node2vec,
CARE-GNN, GREET}.
Following this construction, we combine three symmetric similarity measures for the feature similarity component $\vz^{x}_{uv}$ of the edge $(u,v)\in\cE$ with feature vectors $\vx_u, \vx_v \in \mathbb{R}^{d}$:
\begin{equation}
  \vz^{x}_{uv} \;=\;
  \bigl[\, |\vx_u-\vx_v|,\; \vx_u\odot\vx_v,\; \cos(\vx_u,\vx_v) \,\bigr],
  \label{eq:zfeature}
\end{equation}
where $\odot$ is the Hadamard product, $\cos(\cdot,\cdot)$ the cosine similarity, and
$\vz^{x}_{uv}\in\mathbb{R}^{2d+1}$.
Each component of $\vz^{x}_{uv}$ captures complementary semantic information about the relation between the endpoints, which is useful for characterizing how information relevant to the target label propagates along the edge.
%the absolute difference captures attribute disagreement, the Hadamard product captures co-activation patterns.

The structural block $\vz^{s}_{uv}$ collects local statistics of the two end-points, which include
\begin{equation}
\begin{aligned}
  \vz^{s}_{uv} \;=\;
  \Bigl[\,&
    \log(1{+}d_u)+\log(1{+}d_v),\;
    \bigl|\log\tfrac{1+d_u}{1+d_v}\bigr|,\\
    &J_{uv},\;
    \tfrac{|\cN_u\cap\cN_v|}{\sqrt{d_u d_v}}
  \,\Bigr],
\end{aligned}
  \label{eq:structural}
\end{equation}
with $J_{uv}=|\cN_u\cap\cN_v|/|\cN_u\cup\cN_v|$ the Jaccard overlap of the endpoint neighborhoods $\cN_u$, $\cN_v$ .
The symmetric degree sum and absolute log-ratio encode the magnitude and asymmetry of endpoint connectivity, while the
overlap terms measure triadic closure around the edge.
%These quantities are directly motivated by fraud topology 
Given a camouflaged anomaly node, it can be attached to well-connected normal nodes but rarely shares their neighborhoods, so heterophilic edges concentrate at large degree asymmetry and low overlap~\citep{CARE-GNN,PCGNN}.

The complete edge descriptor $\vz_{uv}$ is constructed as 
\begin{equation}
    \vz_{uv} = \Bigl[\,\vz^{x}_{uv} \;\|\; \vz^{s}_{uv} \Bigr],
    \label{eq:zuv}
\end{equation}
where $\|$ is a concatenation operator.
It is clear that $\vz_{uv}$ preserves the symmetric semantic of the edge in this construction and prevents label leakage of the targeted label into the extracted features.
We employ a shallow MLP network, parameterized by $\phi$, to train a prediction function from the input $\vz_{uv}$ to the homophily/heterophily semantic of the edge.
The training target $t_{uv}$ is based on whether the endpoints of an edge $(u, v)$ belong to the set $\mathcal{V}_L$.
Firstly, the edge set is split according to how many of its endpoints belong to the training set $\mathcal{V}_L$:
\begin{align}
  \mathcal{E}_2 &= \{(u,v)\in\mathcal{E} :   |\{u, v\} \cap \mathcal{V}_L| = 2 \} \label{eq:tuvE2}, \\
  \mathcal{E}_1 &= \{(u,v)\in\mathcal{E} : |\{u, v\} \cap \mathcal{V}_L| = 1 \} \label{eq:tuvE1}
\end{align}
and let $\mathcal{E}_0$ collect the remaining edges.
For $(u,v)\in\mathcal{E}_2\cup\mathcal{E}_1$
we assign the binary heterophily target
\begin{equation}
  t_{uv} =
  \begin{cases}
    \mathds{1}\!\left[\, y_u \neq y_v \,\right], & (u,v)\in\mathcal{E}_2,\\[2pt]
    y_u, & (u,v)\in\mathcal{E}_1,\; u \in \mathcal{V}_L,\\[2pt]
    y_v, & (u,v)\in\mathcal{E}_1,\; v \in \mathcal{V}_L .
  \end{cases}
  \label{eq:edge-target}
\end{equation}
The first case is the definition of edge heterophily and it constructs a clean label edge set for training.
The remaining cases are the same definition evaluated under modal imputation of the missing endpoint.
Because anomalies are rare, the marginally most probable label of an unobserved node is
$\hat{y}=0$, and $\mathds{1}[\,y_u \neq 0\,] = y_u$.
A singly-labeled edge is therefore treated as
heterophilic when its labeled endpoint is anomalous and as homophilic otherwise, which is consistent with the camouflage assumption in fraud graphs.
From the empirical aspect, admitting the approximated targets in $\mathcal{E}_1$ trades label noise against the number of training samples.
Nevertheless, it is expected to have a bias estimate of the heterophily target of the edges in $\mathcal{E}$ under the semi-supervised setting.
We deliberately admit this approximation error rather than restricting training to $\mathcal{E}_2$ because the operator is designed to absorb the noisy estimate.
Proposition~\ref{prop:e1} characterizes exactly what the resulting target estimates, and Corollary~\ref{cor:camouflage} states when the imputation is beneficial.

The trained parameters $\phi$ are then used to produce the predicted soft heterophily score
\begin{equation}
  \hat{q}_{uv} \;=\; \text{MLP}_\phi(\vz_{uv})
  \label{eq:quvfromphi}
\end{equation}
The training loss of Phase 1 is designed as
\begin{equation}
    \mathcal{L}_{EC} = \mathcal{L}_{BCE}(t_{uv}, \hat{q}_{uv}) +   \mathcal{L}_{MSE}(t_{uv}, \hat{q}_{uv})
    \label{eq:LossEC}
\end{equation}
where $\mathcal{L}_{BCE}$ and $\mathcal{L}_{MSE}$ are the binary cross entropy loss and the Brier loss, respectively.
The Brier loss is additionally included to encourage calibrated edge-heterophily probabilities, since edge-level predictions in GNNs can exhibit substantial probability miscalibration \citep{nascimento2024innout}.
To reduce the estimation noise caused by the sparsity of the training data, and to keep the labels of a node out of the scores of its own edges, the calibrator is cross-fitted with folds that are grouped by node.
The labeled nodes $\mathcal{V}_L$ are split into $K$ class-stratified folds $\mathcal{F}_1,\dots,\mathcal{F}_K$, and the $j$-th model $\phi^{(j)}$ is trained on the edges of $\mathcal{E}_2\cup\mathcal{E}_1$ none of whose labeled endpoints lies in $\mathcal{F}_j$.
An edge is then scored by the average of the models whose held-out fold contains one of its endpoints, and by the average of all $K$ models when neither endpoint is labeled.
Proposition~\ref{prop:crossfit} states which scores are free of a node's own label under this rule.
%The full training pipeline is given in Algorithm 1.
After getting the aggregated prediction score $\hat{q}_{uv}$, we set $\hat{q}_{uu} = 0$ i.e. the self-loop edge is always homophilic.
%\begin{algorithm}[h]
%\caption{Edge calibration with K-fold training}
%\label{alg:ecsbwgnn}
%\begin{algorithmic}[1]
%\Require Graph $G$, labeled set $\mathcal{V}_L$, folds $K$
%\State Compute the edge descriptor $z_{uv}$ using Eq.~\eqref{eq:zfeature}
%\State Stratify $\mathcal{V}_L$ into $K$ folds
%\For{$j = 1,\dots,K$}
%  \State Construct the target for fold $j$, based Eq.~\eqref{eq:edge-target}
%  \State Train $MLP_\phi^{(j)}$ on using $\mathcal{L}_{EC}$ in Eq.~\eqref{eq:LossEC}
%\EndFor
%\State Average out-of-fold predictions to get $\hat{q}_{uv}$
%\end{algorithmic}
%\end{algorithm}
